% Reusable dissertation assets
\section{Platform framing}

\KnowledgePlatformVision

\section{Knowledge representation snippets}

\subsection{RDF or Turtle concept graph example}

\begin{lstlisting}[language=Turtle,caption={Illustrative Turtle snippet for the shared concept network.},label={lst:turtle-concept-network}]
@prefix ex: <https://example.org/icg/> .
@prefix skos: <http://www.w3.org/2004/02/skos/core#> .
@prefix dct: <http://purl.org/dc/terms/> .

ex:graph-algorithms a ex:Concept ;
  skos:prefLabel "Graph algorithms" ;
  ex:recommendedFor "middle-school-plus" ;
  ex:hasPrerequisite ex:loops , ex:variables , ex:graphs-intro ;
  ex:supportsService ex:knowledge-atlas , ex:algorithmica .

ex:lesson-plan-graph-paths a ex:LessonPlan ;
  dct:title "Shortest paths and traversals" ;
  ex:selectsConcept ex:graphs-intro , ex:breadth-first-search , ex:dijkstra-intro ;
  ex:deliveryMode ex:course-view .
\end{lstlisting}

\subsection{SPARQL lesson-plan extraction example}

\begin{lstlisting}[language=SPARQL,caption={Representative SPARQL query for extracting an ordered lesson-plan subset.},label={lst:sparql-lesson-plan}]
PREFIX ex: <https://example.org/icg/>
PREFIX skos: <http://www.w3.org/2004/02/skos/core#>

SELECT ?position ?concept ?label
WHERE {
  ex:lesson-plan-graph-paths ex:orderedItem ?item .
  ?item ex:position ?position ;
        ex:concept ?concept .
  ?concept skos:prefLabel ?label .
}
ORDER BY ?position
\end{lstlisting}

\section{Application and service snippets}

\subsection{TypeScript API example}

\begin{lstlisting}[language=TypeScript,caption={Example TypeScript endpoint for retrieving a service-specific lesson plan.},label={lst:typescript-api}]
interface LessonPlanRequest {
  serviceId: "knowledge-atlas" | "algorithmica" | "disciplina";
  conceptIds: string[];
}

export async function buildLessonPlan(request: LessonPlanRequest) {
  const response = await fetch("https://api.example.org/lesson-plans", {
    method: "POST",
    headers: { "Content-Type": "application/json" },
    body: JSON.stringify(request),
  });

  if (!response.ok) {
    throw new Error("Lesson-plan generation failed.");
  }

  return response.json();
}
\end{lstlisting}

\subsection{JSON response example}

\begin{lstlisting}[language=JSON,caption={Example lesson-plan payload returned by the knowledge API.},label={lst:json-lesson-plan}]
{
  "serviceId": "knowledge-atlas",
  "pathTitle": "Intro to graph algorithms",
  "orderedConcepts": [
    { "id": "graphs-intro", "label": "Graphs and nodes" },
    { "id": "breadth-first-search", "label": "Breadth-first search" },
    { "id": "dijkstra-intro", "label": "Shortest-path intuition" }
  ],
  "deliveryMode": "course-view"
}
\end{lstlisting}

\subsection{SQL analytics example}

\begin{lstlisting}[language=SQL,caption={Example analytics table for cross-service usage reporting.},label={lst:sql-analytics}]
CREATE TABLE service_interactions (
  interaction_id INTEGER PRIMARY KEY,
  service_id VARCHAR(32) NOT NULL,
  learner_id VARCHAR(64) NOT NULL,
  concept_id VARCHAR(64) NOT NULL,
  interaction_type VARCHAR(32) NOT NULL,
  completed_at TIMESTAMP NOT NULL
);
\end{lstlisting}

\section{Comparison table template}

\begin{table}[H]
    \centering
    \begin{threeparttable}
        \caption{Example comparison table for client-facing services built on the shared knowledge API.}
        \label{tab:service-comparison}
        \begin{tabularx}{\linewidth}{>{\raggedright\arraybackslash}p{0.18\linewidth} Y Y Y}
            \toprule
            Service & Primary audience & Pedagogical emphasis & Representative outputs \\
            \midrule
            \LandingPageService{} & Prospective clients, schools, and partners & Communication, trust, service discovery & Service overview, case studies, onboarding information \\
            \KnowledgeAtlas{} & Teachers, planners, and course designers & Course-oriented pathway inspection and standards alignment & Lesson plans, pathway views, curriculum maps \\
            \AlgorithmicaApp{} & Learners in a focused topic sequence & Gamified missions and progression through algorithmic concepts & Levels, quests, mastery checkpoints, feedback loops \\
            \DisciplinaApp{} & Learners and facilitators exploring broader concept structures & Guided progression through curated concept subsets & Sequenced challenges, progress views, concept unlocks \\
            \bottomrule
        \end{tabularx}
    \end{threeparttable}
\end{table}

\section{Multi-panel UI screenshot template}

\begin{figure}[p]
    \centering
    \begin{subfigure}[t]{0.48\linewidth}
        \centering
        \ScreenPlaceholder{Replace with a landing-page overview screenshot.}
        \caption{Client-facing landing page}
    \end{subfigure}
    \hfill
    \begin{subfigure}[t]{0.48\linewidth}
        \centering
        \ScreenPlaceholder{Replace with a \KnowledgeAtlas{} course or pathway view.}
        \caption{Course-oriented web application}
    \end{subfigure}

    \vspace{1em}

    \begin{subfigure}[t]{0.48\linewidth}
        \centering
        \ScreenPlaceholder{Replace with an \AlgorithmicaApp{} mission or level-select view.}
        \caption{Gamified mobile experience}
    \end{subfigure}
    \hfill
    \begin{subfigure}[t]{0.48\linewidth}
        \centering
        \ScreenPlaceholder{Replace with a \DisciplinaApp{} concept or progression view.}
        \caption{Alternative learner-facing service}
    \end{subfigure}

    \caption{Reusable multi-panel layout for UI screenshots across the service portfolio.}
    \label{fig:service-ui-grid}
\end{figure}


\section{Evaluation reporting templates}

\subsection{Participant profile template}

\begin{table}[H]
    \centering
    \caption{Template for participant and reviewer reporting.}
    \label{tab:participant-profile-template}
    \begin{tabularx}{\linewidth}{>{\raggedright\arraybackslash}p{0.24\linewidth} Y Y Y}
        \toprule
        Group & Count & Background variables & Notes for analysis \\
        \midrule
        Learners & --- & Age band, prior CS exposure, familiarity with games or graph interfaces & Separate novice and experienced users where relevant \\
        Teachers & --- & Subject area, years of teaching, curriculum responsibilities & Capture both adoption concerns and pedagogical alignment feedback \\
        Expert reviewers & --- & Ontology, HCI, curriculum, or software engineering expertise & Use for heuristic or standards-focused review cycles \\
        \bottomrule
    \end{tabularx}
\end{table}

\subsection{Result summary template}

\begin{table}[H]
    \centering
    \caption{Template for summarizing evaluation outcomes across services.}
    \label{tab:result-summary-template}
    \begin{tabularx}{\linewidth}{>{\raggedright\arraybackslash}p{0.22\linewidth} Y Y Y}
        \toprule
        Dimension & Web app evidence & Mobile app evidence & Interpretation \\
        \midrule
        Navigation and orientation & --- & --- & Explain where graph transparency or guidance helped or failed \\
        Engagement and motivation & --- & --- & Contrast course-oriented behavior with gamified behavior \\
        Learning support & --- & --- & Relate outcomes back to concept sequencing and prerequisite logic \\
        Teacher or client confidence & --- & --- & Report whether the service portfolio looks coherent and trustworthy \\
        \bottomrule
    \end{tabularx}
\end{table}
